\begin{proposition}[I10 candidate: endpoint descent conditional on the preceding whole-family maps]
Let $A,K,B,Q,T$ be as in I9, and assume additionally $T^2=0$.
Fix $i\geq2$. Suppose continuous maps
\[
 \Phi_r:\operatorname{End}_A^r(K)\longrightarrow Q^r,
 \qquad 0\leq r<i,
\]
have already been constructed on the entire endomorphism groups,
with $\Phi_0$ the prescribed diagonal map and $\Phi_1(\theta)=T$.
Assume they respect the diagonal actions and multiplication whenever
the total degree is less than $i$. For primes
$\mathfrak p\subset\mathfrak q$ with height gap $i$, there is then
a unique continuous map
\[
 \Phi_{\mathfrak q\mathfrak p}:
   \operatorname{Hom}_A(E_{\mathfrak p},E_{\mathfrak q})
           \longrightarrow e_{\mathfrak q}Q^i e_{\mathfrak p}
\]
agreeing with the prescribed value of every finite path word between
these endpoints. This asserts continuity on one interval at a time;
it does not assert continuity or convergence of the full family.
\end{proposition}
\begin{proof}
Put $R=A_{\mathfrak q}$, $h=\operatorname{ht}\mathfrak p$, and
$j=\operatorname{ht}\mathfrak q=h+i$. The residual complex $K_R$
is the truncation of $K_A$ to primes contained in $\mathfrak q$;
its last term is $E_{\mathfrak q}$. Write
\[
 C^a=\operatorname{Hom}_R(E_{\mathfrak p},K_R^a).
\]
Every such map is also an $A$-linear map into the corresponding
direct summand of $K_A^a$. The inclusion into the global
endomorphism group is continuous for the finite topology. Define
\[
 \psi:C^{j-1}\longrightarrow e_{\mathfrak q}Q^i e_{\mathfrak p},
 \qquad g\longmapsto e_{\mathfrak q}T\Phi_{i-1}(g).
                                                        \tag{I10.1}
\]
This is continuous because the \emph{entire} preceding-degree map
is assumed continuous. Let $a_{\mathfrak q}\in B$ be the diagonal
idempotent selecting primes contained in $\mathfrak q$.
For $b\in C^{j-2}$, its local differential is
$a_{\mathfrak q}\theta b$. The inductive multiplication hypothesis
therefore gives
\[
 \psi(db)
  =e_{\mathfrak q}T a_{\mathfrak q}T\Phi_{i-2}(b)=0.     \tag{I10.2}
\]
Indeed the incidence relation, first on finite diagonal truncations
and then by continuity, says that every source of an arrow ending
at $\mathfrak q$ is already selected by $a_{\mathfrak q}$.
Thus the last expression equals
$e_{\mathfrak q}T^2\Phi_{i-2}(b)$.

Since $j-1>h$, I6 identifies $\ker(d:C^{j-1}\to C^j)$ with
the closure of $d(C^{j-2})$. Hence (I10.2) and continuity imply
that $\psi$ kills this kernel. I7b says that
\[
 C^{j-1}\longrightarrow\operatorname{im}(d)\subset C^j
\]
is an open quotient because $j>h+1$. Thus (I10.1) descends to a
continuous map on $\operatorname{im}(d)$. I6 makes this image dense
in all of $C^j$, because $K_R$ has no term above degree $j$.
The target corner is a closed subgroup of the complete group $Q^i$,
being the image of the continuous idempotent projector
$x\mapsto e_{\mathfrak q}xe_{\mathfrak p}$.
Uniform continuity of additive homomorphisms therefore extends the
map uniquely to $C^j$.

The $R$-linear and $A$-linear Hom groups here coincide because
elements outside $\mathfrak q$ act invertibly on both modules.
This gives the displayed map in the proposition. It is linear for
the left $A_{\mathfrak q}$-action on the dense image and for the
right $T_{\mathfrak p}$-action. Approximating an element of
$\widehat{A_{\mathfrak q}}$ by elements of $A_{\mathfrak q}$,
and using continuity on the endpoint Hom and in $Q$, proves the
full completed left-diagonal compatibility. A path with final arrow
$\delta_{\mathfrak r\mathfrak q}$ and preceding map supported
at $\mathfrak r$ is the local differential of that preceding map.
Formula (I10.1), the inductive path formula, and the just-proved
diagonal compatibility therefore give every finite path its
prescribed value.
I5 supplies density of these path words on the endpoint Hom group,
which also proves uniqueness.
\end{proof}

\paragraph{Scope of the remaining induction.}
If I9 is accepted, I10 yields continuous maps for every height-gap-two
interval. It does \emph{not} automatically supply the whole degree-two
map required to repeat the proposition for height gap three. Finite
topology on $\operatorname{Hom}(K^h,K^j)$ is a product over source
columns, but the target in (I10.1) must already admit continuous
assembly of the whole preceding-degree family. Establishing that
assembly in each new degree, or constructing a countertarget to it,
remains the precise broad-topology gate.
